\bookchapter{Preface}{ch:00-preface}

\begin{chapterintro}

This sample book exists to exercise every feature of the series template in a few pages. It explains, briefly and truthfully, what a text file is and how to count the words in one with Python. Read it next to its Markdown source to see which convention produces which result.

\end{chapterintro}

\section{Who This Book Is For}\label{00-preface:who-this-book-is-for}

You can open a terminal and type a command. You do \textbf{not} need to know anything about files, bytes, or Python beyond that. Everything else is explained when it's needed.

\section{How This Book Is Organized}\label{00-preface:how-this-book-is-organized}

\begin{itemize}
\tightlist
\item
  \textbf{Part I} builds the mental model: what is actually inside a text file.
\item
  \textbf{Part II} puts it to use: a ten-line Python program that counts words, and a lab.
\item
  \textbf{Appendix A} is a cheat sheet; \textbf{Appendix B} is the glossary.
\end{itemize}

\section{How to Use This Book}\label{00-preface:how-to-use-this-book}

Read the chapters in order; they're short. \textbf{Type the code yourself} rather than copying it. Every chapter ends with a lab and a \textbf{Summary, Key Terms, and Review Questions} section.

\section{What You Need Installed}\label{00-preface:what-you-need-installed}

\begin{itemize}
\tightlist
\item
  \textbf{Python 3.12 or newer} (check with \texttt{python3\ -\/-version})
\item
  \textbf{A terminal}: Terminal on macOS, any terminal on Linux, WSL on Windows
\end{itemize}

\section{Conventions Used in This Book}\label{00-preface:conventions-used-in-this-book}

\begin{itemize}
\tightlist
\item
  \textbf{Bold} marks a new term the first time it is explained. Every such term is also in Appendix B.
\item
  \texttt{Constant\ width} marks code, commands, file names, and anything else you would type exactly.
\item
  Code blocks show Python, shell commands, or program output. Shell commands are shown without a prompt, so you can paste them as they are.
\item
  Three kinds of callout appear in the text:
\end{itemize}

\begin{callout}{tip}

A suggestion that saves time or trouble.

\end{callout}

\begin{callout}{note}

A general remark, or something to be aware of.

\end{callout}

\begin{callout}{warning}

A mistake that can cost you data, money, or security. Read these twice.

\end{callout}
